\section{Literature Review: Key Recent Publications}
\label{sec:literature_review}

The following section provides structured academic summaries of five significant recent publications that together define the state of the art in multi-omics research in LBD, spanning genomics, epigenomics, metabolomics, single-cell immunology, and translational biomarker development.

\subsection{Publication 1: Large-Scale GWAS and Multi-Omics Integration (2025)}
\textbf{Citation:} International LBD Genomics Consortium, ``Genome-wide association study provides insights into the genetic architecture of Lewy body dementia'' \cite{chia2025gwas}.

\begin{itemize}
    \item \textbf{Research Problem:} Despite previous GWAS efforts \cite{chia2021lbd}, the genetic architecture of LBD remained incompletely characterised due to limited sample sizes. A critical gap was the absence of functionally resolved multi-omics integration connecting genetic risk variants to downstream biological effects.
    \item \textbf{Methodology:} A meta-analysis of genome-wide association data aggregating 4,252 LBD cases and 189,290 controls, drawn from multiple international cohorts. Multi-omics integration was performed using expression quantitative trait loci (eQTL) mapping, Mendelian randomisation, and enrichment analyses across brain cell type-specific regulatory data.
    \item \textbf{Findings:} Confirmed known risk loci (\textit{APOE}, \textit{GBA}, \textit{BIN1}, \textit{SNCA-AS1}) and identified a novel genome-wide significant locus at \textit{SYT16} (encoding synaptotagmin-16, a regulator of synaptic vesicle fusion). Multi-omics integration pinpointed 85 LBD risk genes and 51 enriched biological pathways, with several gene-expression changes observed in LBD post-mortem brain tissues. Potential druggable targets within the identified pathway network were highlighted.
    \item \textbf{Strengths:} The largest LBD genetic study to date; the use of multi-omics integration to move from risk loci to risk genes is a methodological step-change. The functional annotation across cell types (neurons, oligodendrocytes, microglia) adds biological granularity.
    \item \textbf{Limitations:} GWAS identifies common variants; rare pathogenic variants in \textit{GBA1} and \textit{SNCA} are not well captured. The multi-omics integration used post-mortem brain data that may not reflect early or blood-accessible changes. Cross-ethnic representation remains limited.
    \item \textbf{Relevance to Proposed Project:} Provides the definitive genetic foundation for this project. The 85 LBD risk genes constitute a high-priority candidate list for downstream omics prioritisation. The identified pathways (synaptic vesicle trafficking, lysosomal function, neuroinflammation) will anchor the biological interpretation layer of our integration pipeline.
\end{itemize}

\subsection{Publication 2: DNA Methylation Meta-Analysis in LBD Brain Tissue (2024)}
\textbf{Citation:} ``Interrogating DNA methylation associated with Lewy body pathology in a large multi-cohort meta-analysis'' \cite{nabais2024methylation}.

\begin{itemize}
    \item \textbf{Research Problem:} DNA methylation (DNAm) is a central mechanism by which genetic risk variants exert epigenetic effects on gene expression. Prior DNAm studies in LBD were underpowered, conducted in single brain regions, and lacked cross-cohort validation, resulting in limited reproducibility.
    \item \textbf{Methodology:} A large-scale epigenome-wide association study (EWAS) meta-analysis using 1,239 cortical brain tissue samples from multiple international cohorts. DNAm profiles (EPIC/450K arrays) were harmonised, batch-corrected, and analysed for differential methylation in relation to Lewy body pathological burden, controlling for age, sex, cell-type proportions, and genetic ancestry.
    \item \textbf{Findings:} Thirty genome-wide significant differentially methylated loci were identified. Enrichment was observed in genes associated with neuroinflammation (\textit{UBASH3B}), platelet-activating factor receptor (\textit{PTAFR}), and notably within the \textit{SNCA} genomic region---directly connecting the principal genetic risk gene to epigenetic dysregulation. Pathway enrichment revealed functional clusters related to synaptic function, vascular alterations, and immune activation.
    \item \textbf{Strengths:} The largest brain-based EWAS in LBD to date; multi-region sampling and multi-cohort design substantially improve reproducibility compared to prior studies. The integration of genetic co-localisation analyses strengthens causal inference.
    \item \textbf{Limitations:} Restricted to post-mortem brain tissue, which limits direct biomarker translation. DNAm in brain tissue does not perfectly reflect peripheral blood methylation patterns. End-stage pathology may obscure changes present in earlier disease phases.
    \item \textbf{Relevance to Proposed Project:} The 30 identified loci provide a high-confidence epigenomic layer for integration with genomic and metabolomic data. The \textit{SNCA} locus methylation findings are particularly compelling for cross-layer analysis, linking genetics, epigenetics, and potential expression changes in a single gene of central pathological importance.
\end{itemize}

\subsection{Publication 3: Plasma Metabolomics Profiling in DLB (2024)}
\textbf{Citation:} ``Translational advances in Lewy body diseases: plasma metabolomics profiling distinguishes dementia with Lewy bodies from Alzheimer's disease'' \cite{frontiers2024metabolomics}.

\begin{itemize}
    \item \textbf{Research Problem:} Reliable differential diagnosis between DLB and AD is a major clinical challenge due to overlapping symptoms and the lack of cost-effective biomarkers. Metabolomics offers the potential for blood-accessible, quantitative molecular profiles that reflect disease-specific biochemical states.
    \item \textbf{Methodology:} Untargeted plasma metabolomics profiling (630 metabolites, >200 derived metabolism indicators) was applied to cohorts of clinically diagnosed DLB patients, AD patients, and age-matched healthy controls. Multivariate statistical methods (principal component analysis, partial least squares discriminant analysis) were used for group separation, followed by pathway enrichment and metabolite set enrichment analyses.
    \item \textbf{Findings:} Significant separation of DLB from both AD and healthy controls was achieved. Disease-specific alterations were identified in neurotransmitter metabolism (serotonin, dopamine pathways), branched-chain amino acids, and sphingolipid/ceramide metabolism---the latter being of direct relevance to \textit{GBA1} function and lysosomal biology. Several metabolite features were proposed as candidate differential diagnostic markers.
    \item \textbf{Strengths:} Blood-accessible profiling is directly translatable to clinical settings. The simultaneous comparison against AD controls is essential for a clinically useful biomarker panel. The breadth of the metabolite coverage (630 features) allows comprehensive pathway reconstruction.
    \item \textbf{Limitations:} Clinical diagnosis without neuropathological confirmation introduces misclassification risk. Cross-platform metabolomics comparisons are difficult due to lack of assay standardisation. The relatively small cohort sizes may limit statistical power for individual metabolite associations.
    \item \textbf{Relevance to Proposed Project:} Demonstrates proof-of-concept that blood metabolomics can differentiate DLB from AD, directly supporting the biomarker objective of this proposal. The sphingolipid and ceramide findings provide a biological bridge between the \textit{GBA1} genetic risk layer and metabolic dysfunction, creating a hypothesis for multi-omics integration.
\end{itemize}

\subsection{Publication 4: Single-Cell Peripheral Immunoprofiling in LBD (2024)}
\textbf{Citation:} ``Single-cell peripheral immunoprofiling of Lewy body disease in a multi-site cohort reveals disease-specific immune signatures'' \cite{singlecell2024immune}.

\begin{itemize}
    \item \textbf{Research Problem:} Neuroinflammation is increasingly recognised as a key contributor to LBD pathophysiology, yet its systemic immune correlates in peripheral blood have been poorly characterised at single-cell resolution. Bulk transcriptomics masks cell-type-specific signals that single-cell RNA sequencing (scRNA-seq) can resolve.
    \item \textbf{Methodology:} Single-cell profiling (scRNA-seq and/or CyTOF mass cytometry) was applied to peripheral blood mononuclear cells (PBMCs) from LBD patients, Parkinson's disease patients without dementia, and healthy controls in a multi-site cohort design. Cell type proportions, differential gene expression, and signalling pathway activities were characterised for each immune cell subset.
    \item \textbf{Findings:} LBD patients showed distinct peripheral immune signatures not observed in PD without dementia or healthy controls, including altered monocyte activation states, changes in natural killer (NK) cell proportions, and unique signalling pathway activities in T-cell subsets. These systemic immune alterations suggest that peripheral immune dysregulation may contribute to, or reflect, central neuroinflammation in LBD.
    \item \textbf{Strengths:} Single-cell resolution eliminates the confounding effect of cell-type composition changes in bulk analyses. Multi-site design strengthens reproducibility. Peripheral blood profiling is minimally invasive and clinically translatable.
    \item \textbf{Limitations:} The relationship between peripheral immune changes and central neuroinflammation remains correlative without mechanistic follow-up. Single-cell data is computationally intensive and requires careful batch correction across sites. Longitudinal data are needed to establish whether immune changes precede or follow onset of dementia.
    \item \textbf{Relevance to Proposed Project:} Provides a cellular-resolution immunomics layer that can be integrated with genomic risk variants implicated in immune function (\textit{BIN1}, \textit{TMEM175}) and epigenetic changes in neuroinflammation-related genes (\textit{UBASH3B}). The multi-site design also provides a valuable model for how to structure this project's own data harmonisation strategy.
\end{itemize}

\subsection{Publication 5: Comprehensive Omics Biomarker Review in Synucleinopathies (2025)}
\textbf{Citation:} ``Biomarkers for Lewy body diseases and other alpha-synucleinopathies: omics-based approaches from biofluids to multi-marker panels'' \cite{boiten2025biomarkers}.

\begin{itemize}
    \item \textbf{Research Problem:} The biomarker field for LBD and related synucleinopathies has expanded rapidly, but evidence is fragmented across platforms, biofluids, and disease stages. A systematic synthesis is needed to identify which omics modalities and biofluid sources offer the most clinically actionable signal.
    \item \textbf{Methodology:} Comprehensive narrative and systematic review of omics-based biomarker studies in synucleinopathies. Studies using CSF, blood (plasma/serum), urine, and olfactory mucosa as source biofluids were included. Biomarker categories reviewed encompassed proteins (alpha-synuclein, NfL, GFAP, pTau181), metabolites, lipids, extracellular vesicle cargo, and epigenetic markers.
    \item \textbf{Findings:} Neurofilament light chain (NfL) and glial fibrillary acidic protein (GFAP) are currently the most analytically mature blood-based markers for neurodegeneration staging in LBD, though they lack disease specificity. Alpha-synuclein SAAs in CSF show high sensitivity and specificity for synucleinopathies but do not sub-classify LBD. Multi-marker panels combining protein, metabolite, and epigenetic features are proposed as the highest-yield future direction, though prospective validation data remain scarce.
    \item \textbf{Strengths:} Broad coverage across biofluids and omics modalities. Practical commentary on analytical platforms and clinical translatability distinguishes this review from purely academic summaries. Positions multi-marker integration as the next frontier in the field.
    \item \textbf{Limitations:} As a review, it does not generate new primary data. The field is rapidly evolving and some cited studies may have since been refined or contradicted. Standardisation of pre-analytical variables (sample collection, storage) across studies remains a critical unresolved challenge.
    \item \textbf{Relevance to Proposed Project:} Provides a state-of-the-art map of which biomarkers already have analytical validation and which gaps this project can address. The explicit call for integrated multi-marker panels directly validates the rationale of this proposal and identifies the current evidence ceiling that we aim to raise.
\end{itemize}

\subsection{Synthesis: The Case for Multi-Omics Integration}
The five studies reviewed above collectively reveal a consistent pattern: each individual omics layer---genomics, epigenomics, metabolomics, single-cell immunomics, or proteomics---provides partial but complementary insight into LBD pathophysiology. No single layer is sufficient for accurate early diagnosis, reliable subtype classification, or confident pathway identification. The convergence of risk genes (\textit{GBA1}, \textit{SNCA}, \textit{BIN1}) in multiple omics layers (GWAS, EWAS, single-cell immune, metabolomics) strongly suggests that integration of these layers will yield molecular signatures of greater sensitivity, specificity, and mechanistic interpretability than any layer alone. This convergence forms the central scientific rationale of the proposed project.
